\subsection{Triangularization over Complex and Real Fields}

\begin{definition} \label{def:diagonal-and-upper-triangular-matrices} \hlt{Diagonal and Upper-Triangular Matrices}\\
The \hlt{diagonal} of a square matrix is its upper-left to lower-right entries. A square matrix is \hlt{diagonal} if all off-diagonal entries are zero, and \hlt{upper triangular} if all entries below the diagonal are zero.
\end{definition}

\begin{theorem} \label{thm:complex-operator-has-eigenvalue} \hlt{Complex Operators Have Eigenvalues}\\
Every operator on a finite-dimensional, nonzero complex vector space has an eigenvalue.
\begin{proof}
Let $T\in\LL(\mathcal{V})$. By Proposition \ref{prop:finite-dimensional-operators-have-annihilating-polynomials}, $T$ has a nonconstant annihilating polynomial. Since $\mathcal{V}$ is complex and nonzero, Theorem \ref{thm:annihilating-polynomials-control-spectrum} implies that one root of this polynomial is an eigenvalue of $T$.\\
\hlt{Structural Commentary}: Finite dimension gives an annihilating polynomial; over $\mathbb{C}$, it factors linearly.\\
\hlt{Key Insight}: One noninjective factor produces an eigenvector.
\end{proof}
\end{theorem}

\begin{theorem} \label{thm:upper-triangular-criteria} \hlt{Criteria for an Upper-Triangular Operator Matrix}\\
Let $T\in\LL(\mathcal{V})$ and let $\beta=(\mathbf{v}_1,\ldots,\mathbf{v}_n)$ be a basis of $\mathcal{V}$. The following are equivalent:
\begin{enumerate}[label=(\roman*)]
\setlength{\itemsep}{0pt}
\item $[T]_{\beta\leftarrow\beta}$ is upper triangular;
\item $T\mathbf{v}_j\in\operatorname{span}(\mathbf{v}_1,\ldots,\mathbf{v}_j)$ for each $j=1,\ldots,n$;
\item $\operatorname{span}(\mathbf{v}_1,\ldots,\mathbf{v}_j)$ is invariant under $T$ for each $j=1,\ldots,n$.
\end{enumerate}
\begin{proof}
\begin{enumerate}[label=(\roman*).]
\setlength{\itemsep}{0pt}
\item[$(i)$]$\Rightarrow (ii)$
The $j$th column of $[T]_{\beta\leftarrow\beta}$ is the coordinate column of $T\mathbf{v}_j$ in the basis $\beta$. Since the matrix is upper triangular, this column has zero entries below position $j$, so $T\mathbf{v}_j\in\operatorname{span}(\mathbf{v}_1,\ldots,\mathbf{v}_j)$.
\item[$(ii)$]$\Rightarrow (i)$
If $T\mathbf{v}_j\in\operatorname{span}(\mathbf{v}_1,\ldots,\mathbf{v}_j)$ for every $j$, then the $j$th coordinate column of $T\mathbf{v}_j$ has zero entries below position $j$. Hence $[T]_{\beta\leftarrow\beta}$ is upper triangular.
\item[$(ii)$]$\Rightarrow (iii)$
If $\mathbf{u}\in\operatorname{span}(\mathbf{v}_1,\ldots,\mathbf{v}_j)$, then linearity gives $T\mathbf{u}$ as a linear combination of $T\mathbf{v}_1,\ldots,T\mathbf{v}_j$, all of which lie in $\operatorname{span}(\mathbf{v}_1,\ldots,\mathbf{v}_j)$. Hence this span is invariant.
\item[$(iii)$]$\Rightarrow (ii)$
Apply invariance of $\operatorname{span}(\mathbf{v}_1,\ldots,\mathbf{v}_j)$ to $\mathbf{v}_j$.
\end{enumerate}
\hlt{Structural Commentary}: Upper triangularity records an invariant flag.\\
\hlt{Key Insight}: Columns below the diagonal vanish exactly when $T\mathbf{v}_j$ uses only earlier basis vectors and itself.
\end{proof}
\end{theorem}

\begin{corollary} \label{cor:upper-triangular-orthonormal-basis} \hlt{Upper-Triangular Form in an Orthonormal Basis}\\
Let $\mathcal{V}$ be a finite-dimensional inner-product space and let $T\in\LL(\mathcal{V})$. If $T$ has an upper-triangular matrix with respect to some basis of $\mathcal{V}$, then $T$ has an upper-triangular matrix with respect to some orthonormal basis of $\mathcal{V}$.
\begin{proof}
Let $\beta=(\mathbf{v}_1,\ldots,\mathbf{v}_n)$ be a basis for which $[T]_{\beta\leftarrow\beta}$ is upper triangular. By Theorem \ref{thm:upper-triangular-criteria}, each subspace
\begin{align}
\mathcal{U}_j=\operatorname{span}(\mathbf{v}_1,\ldots,\mathbf{v}_j),\qquad j=1,\ldots,n,
\nonumber
\end{align}
is invariant under $T$. Apply Gram--Schmidt to $\beta$ and obtain an orthonormal basis $\gamma=(\mathbf{e}_1,\ldots,\mathbf{e}_n)$ with
\begin{align}
\operatorname{span}(\mathbf{e}_1,\ldots,\mathbf{e}_j)=\operatorname{span}(\mathbf{v}_1,\ldots,\mathbf{v}_j)=\mathcal{U}_j.
\nonumber
\end{align}
Thus each $\operatorname{span}(\mathbf{e}_1,\ldots,\mathbf{e}_j)$ is invariant under $T$, and Theorem \ref{thm:upper-triangular-criteria} shows that $[T]_{\gamma\leftarrow\gamma}$ is upper triangular.\\
\hlt{Structural Commentary}: Upper triangularity is preservation of a nested flag of subspaces. Gram--Schmidt preserves that flag while making the basis orthonormal.\\
\hlt{Key Insight}: Triangular form depends on the flag, not on the particular basis used inside the flag.
\end{proof}
\end{corollary}

\begin{theorem} \label{thm:complex-upper-triangularization} \hlt{Complex Upper Triangularisation}\\
Let $\mathcal{V}$ be a finite-dimensional complex vector space and let $T\in\LL(\mathcal{V})$. Then $T$ has an upper-triangular matrix with respect to some basis of $\mathcal{V}$.
\begin{proof} Induct on $n=\dim(\mathcal{V})$. If $n\leq1$, the result is immediate.\\
Suppose $n>1$ and the result holds in smaller dimensions. Choose an eigenvector $\mathbf{v}_1$ of $T$ by Theorem \ref{thm:complex-operator-has-eigenvalue}, and set $\mathcal{U}=\operatorname{span}(\mathbf{v}_1)$. Then $\mathcal{U}$ is $T$-invariant, so the quotient operator $T/\mathcal{U}$ is defined on $\mathcal{V}/\mathcal{U}$, which has dimension $n-1$. By induction, choose a basis $(\mathbf{v}_2+\mathcal{U},\ldots,\mathbf{v}_n+\mathcal{U})$ of $\mathcal{V}/\mathcal{U}$ with respect to which $T/\mathcal{U}$ is upper triangular, and choose representatives $\mathbf{v}_2,\ldots,\mathbf{v}_n\in\mathcal{V}$.\\
The lifted list $(\mathbf{v}_1,\ldots,\mathbf{v}_n)$ is a basis of $\mathcal{V}$: any relation descends to the quotient basis, forcing all coefficients of $\mathbf{v}_2,\ldots,\mathbf{v}_n$ to be zero, and then the coefficient of $\mathbf{v}_1$ is zero. Now $T\mathbf{v}_1\in\operatorname{span}(\mathbf{v}_1)$.\\
For $j\geq2$, upper triangularity of $T/\mathcal{U}$ gives $T\mathbf{v}_j+\mathcal{U} =(T/\mathcal{U})(\mathbf{v}_j+\mathcal{U}) \in\operatorname{span}(\mathbf{v}_2+\mathcal{U},\ldots,\mathbf{v}_j+\mathcal{U})$, so $T\mathbf{v}_j\in\operatorname{span}(\mathbf{v}_1,\ldots,\mathbf{v}_j)$. Hence Theorem \ref{thm:upper-triangular-criteria} gives an upper-triangular matrix for $T$ in the basis $(\mathbf{v}_1,\ldots,\mathbf{v}_n)$.\\
\hlt{Structural Commentary}: Lift an upper-triangular quotient basis above an eigenline.\\
\hlt{Key Insight}: The quotient builds the remaining invariant flag.
\end{proof}
\end{theorem}

\begin{definition} \label{def:complexification-real-vector-space} \label{def:complexification-real-imaginary-parts} \hlt{Complexification of a Real Vector Space}\\
Let $\mathcal{V}$ be a real vector space.
\begin{enumerate}[label=(\roman*)]
\setlength{\itemsep}{0pt}
\item The \hlt{complexification} of $\mathcal{V}$ is $\mathcal{V}_{\mathbb{C}}\defeq\mathcal{V}\times\mathcal{V}$, with $(\mathbf{u},\mathbf{v})$ denoted $\mathbf{u}+i\mathbf{v}$.
\item Addition and $\mathbb{C}$-scalar multiplication are
\begin{align}
(\mathbf{u}_1+i\mathbf{v}_1)+(\mathbf{u}_2+i\mathbf{v}_2)
&=(\mathbf{u}_1+\mathbf{u}_2)+i(\mathbf{v}_1+\mathbf{v}_2), \nonumber\\
(a+bi)(\mathbf{u}+i\mathbf{v})
&=(a\mathbf{u}-b\mathbf{v})+i(b\mathbf{u}+a\mathbf{v}),
\qquad a,b\in\mathbb{R}. \nonumber
\end{align}
\item \hlt{Complexification map} is real-linear embedding $\operatorname{cpx}:\mathcal{V}\hookrightarrow\mathcal{V}_{\mathbb{C}}$; $\operatorname{cpx}(\mathbf{u})=\mathbf{u}+i\mathbf{0}$, where $\{\mathbf{u}+i\mathbf{0}:\mathbf{u}\in\mathcal{V}\}$.
\item For $\mathbf{z}=\mathbf{u}+i\mathbf{v}\in\mathcal{V}_{\mathbb{C}}$, define $\operatorname{Re}\mathbf{z}\defeq\mathbf{u}$ and $\operatorname{Im}\mathbf{z}\defeq\mathbf{v}$. Hence $\mathbf{z}=\mathbf{0}$ iff $\operatorname{Re}\mathbf{z}=\operatorname{Im}\mathbf{z}=\mathbf{0}$.
\end{enumerate}
\end{definition}

\begin{proposition} \label{prop:complexification-vector-space-basis} \hlt{Complexification Preserves Bases}\\
Let $\mathcal{V}$ be real. Then $\mathcal{V}_{\mathbb{C}}$ is complex, and every real basis $(\mathbf{v}_1,\ldots,\mathbf{v}_n)$ of $\mathcal{V}$ is a complex basis of $\mathcal{V}_{\mathbb{C}}$.\\
Thus $\dim_{\mathbb{C}}(\mathcal{V}_{\mathbb{C}})=\dim_{\mathbb{R}}(\mathcal{V})$.
\begin{proof}
The vector-space axioms follow componentwise from those of $\mathcal{V}$.\\
For $\mathbf{u},\mathbf{v}\in\mathcal{V}$, write $\mathbf{u}=\sum\limits_{j=1}^{n}a_j\mathbf{v}_j$ and $\mathbf{v}=\sum\limits_{j=1}^{n}b_j\mathbf{v}_j$, $a_j,b_j\in\mathbb{R}$. Then $\mathbf{u}+i\mathbf{v}=\sum\limits_{j=1}^{n}(a_j+ib_j)\mathbf{v}_j$, so the list spans.\\
If $\sum\limits_{j=1}^{n}\lambda_j\mathbf{v}_j=\mathbf{0}$ and $\lambda_j=a_j+ib_j$, then $\sum\limits_{j=1}^{n}a_j\mathbf{v}_j=\mathbf{0}$ and $\sum\limits_{j=1}^{n}b_j\mathbf{v}_j=\mathbf{0}$.\\
Real linear independence gives $a_j=b_j=0$ for all $j$, hence $\lambda_j=0$ for all $j$.\\
\hlt{Structural Commentary}: Complexification keeps the basis length unchanged.\\
\hlt{Key Insight}: Separate real and imaginary parts.
\end{proof}
\end{proposition}

\begin{definition} \label{def:complexification-real-operator} \hlt{Complexification of a Real Operator}\\
Let $\mathcal{V}$ be a real vector space and let $T\in\LL(\mathcal{V})$.\\
The \hlt{complexification} of $T$ is the operator $T_{\mathbb{C}}\in\LL(\mathcal{V}_{\mathbb{C}})$ defined by $T_{\mathbb{C}}(\mathbf{u}+i\mathbf{v})\defeq T\mathbf{u}+iT\mathbf{v}$.
\end{definition}

\begin{proposition} \label{prop:complexification-operator-basic-properties} \hlt{Complexification Preserves Operator Data}\\
Let $\mathcal{V}$ be a real vector space and let $T\in\LL(\mathcal{V})$.
\begin{enumerate}[label=(\roman*)]
\setlength{\itemsep}{0pt}
\item $T_{\mathbb{C}}\in\LL(\mathcal{V}_{\mathbb{C}})$.
\item $(T_{\mathbb{C}})^k(\mathbf{u}+i\mathbf{v})=T^k\mathbf{u}+iT^k\mathbf{v}$ for every integer $k\geq0$.
\item If $\beta$ is a real basis of $\mathcal{V}$, then $[T_{\mathbb{C}}]_{\beta\leftarrow\beta}=[T]_{\beta\leftarrow\beta}$, now regarded as a complex matrix.
\end{enumerate}
\begin{proof}
\begin{enumerate}[label=(\roman*)]
\setlength{\itemsep}{0pt}
\item Additivity is immediate. For $\alpha=a+bi\in\mathbb{C}$,
\begin{align}
T_{\mathbb{C}}\bigl(\alpha(\mathbf{u}+i\mathbf{v})\bigr) =T(a\mathbf{u}-b\mathbf{v})+iT(b\mathbf{u}+a\mathbf{v}) =(aT\mathbf{u}-bT\mathbf{v})+i(bT\mathbf{u}+aT\mathbf{v}) =\alpha T_{\mathbb{C}}(\mathbf{u}+i\mathbf{v}),
\nonumber
\end{align}
so $T_{\mathbb{C}}$ is complex-linear.
\item For $k=0$, the claim is the identity. If it holds for $k$, then
\begin{align}
(T_{\mathbb{C}})^{k+1}(\mathbf{u}+i\mathbf{v}) =T_{\mathbb{C}}(T^k\mathbf{u}+iT^k\mathbf{v}) =T^{k+1}\mathbf{u}+iT^{k+1}\mathbf{v}.
\nonumber
\end{align}
Thus the formula holds for all $k\geq0$.
\item If $T\mathbf{v}_j=\sum\limits_i a_{ij}\mathbf{v}_i$ with $a_{ij}\in\mathbb{R}$, then $T_{\mathbb{C}}\mathbf{v}_j=\sum\limits_i a_{ij}\mathbf{v}_i$. Hence the coordinate matrix is unchanged.
\end{enumerate}
\hlt{Structural Commentary}: Complexification keeps the coordinate matrix unchanged.\\
\hlt{Key Insight}: Same real entries; larger scalar field.
\end{proof}
\end{proposition}

\begin{figure}[H]
  \centering
  \begin{tikzpicture}[
  >=Stealth,
  every node/.style={font=\small},
  box/.style={draw=gray!55, rounded corners=3pt, inner xsep=7pt, inner ysep=4.5pt, align=center},
  arrow/.style={->, line width=0.85pt, blue!70!black}
]
  \node[box] (real) at (0,0) {$\mathbf{u},\mathbf{v}\in\mathcal{V}$};
  \node[box] (complex) at (2.65,0) {$\mathbf{u}+i\mathbf{v}\in\mathcal{V}_{\mathbb{C}}$};
  \node[box] (image) at (5.80,0) {$T\mathbf{u}+iT\mathbf{v}$};

  \draw[arrow] (real) -- (complex);
  \draw[arrow] (complex) -- (image);
\end{tikzpicture}

  \caption{Complexification packages two real vectors: $T_{\mathbb{C}}(\mathbf{u}+i\mathbf{v})=T\mathbf{u}+iT\mathbf{v}$.}
  \label{fig:complexification-real-operator}
\end{figure}

\begin{definition} \label{def:almost-upper-triangular-real-matrix} \hlt{Almost Upper-Triangular Real Matrices}\\
Let $A\in\operatorname{Mat}(n,n,\mathbb{R})$. The matrix $A$ is \hlt{almost upper triangular} if there exist $d_1,\ldots,d_m\in\{1,2\}$ with $\sum_jd_j=n$ such that, in the corresponding block form $A=[A_{ij}]_{i,j=1}^{m}$,
\begin{enumerate}[label=(\roman*)]
\setlength{\itemsep}{0pt}
\item $A_{ij}=0$ for $i>j$;
\item each diagonal block is either $A_{jj}=[\lambda_j]$, $\lambda_j\in\mathbb{R}$, or $A_{jj}=
\begin{bmatrix}
a_j&-b_j\\
b_j&a_j
\end{bmatrix}$, $a_j,b_j\in\mathbb{R},\ b_j\neq0$.
\end{enumerate}
The matrix $A$ is in \hlt{real Schur block form} if only (i) holds for some such block sizes $d_j\in\{1,2\}$.
\end{definition}

\begin{theorem} \label{thm:real-almost-upper-triangularization} \hlt{Real Almost Upper Triangularisation}\\
Let $\dim_{\mathbb{R}}\mathcal{V}<\infty$ and let $T\in\LL(\mathcal{V})$. Then $\exists$ a basis $\beta$ of $\mathcal{V}$ such that $[T]_{\beta\leftarrow\beta}$ is almost upper triangular.
\begin{proof}
Induct on $n=\dim_{\mathbb{R}}\mathcal{V}$. If $n\leq1$, clear.\\
For $n\geq2$, Theorem \ref{thm:complex-operator-has-eigenvalue} gives an eigenpair $T_{\mathbb{C}}(\mathbf{u}+i\mathbf{v})=(a+bi)(\mathbf{u}+i\mathbf{v})$, where $\mathbf{u}+i\mathbf{v}\neq\mathbf{0}$.\\
Thus $T\mathbf{u}=a\mathbf{u}-b\mathbf{v}$ and $T\mathbf{v}=b\mathbf{u}+a\mathbf{v}$.\\
If $b=0$, choose $\mathbf{0}\neq\mathbf{x}\in\{\mathbf{u},\mathbf{v}\}$, set $\mathcal{U}=\operatorname{span}(\mathbf{x})$, $\alpha=(\mathbf{x})$, so $[T|_{\mathcal{U}}]_{\alpha\leftarrow\alpha}=[a]$.\\
If $b\neq0$, set $\mathcal{U}=\operatorname{span}(\mathbf{u},\mathbf{v})$, $\alpha=(\mathbf{u},-\mathbf{v})$; then $\mathbf{u},\mathbf{v}$ are independent and $[T|_{\mathcal{U}}]_{\alpha\leftarrow\alpha}=
\begin{bmatrix}
a&-b\\
b&a
\end{bmatrix}$.\\
In either case, $T(\mathcal{U})\subseteq\mathcal{U}$, $\dim\mathcal{U}\in\{1,2\}$, and the first diagonal block is admissible.\\
By induction, $T/\mathcal{U}$ has an almost upper-triangular matrix in some quotient basis $\overline{\gamma}$.\\
Lift $\overline{\gamma}$ to $\gamma\subseteq\mathcal{V}$ and set $\beta=(\alpha,\gamma)$. Then $[T]_{\beta\leftarrow\beta}$ is block upper triangular with first block $[T|_{\mathcal{U}}]_{\alpha\leftarrow\alpha}$ and quotient block $[T/\mathcal{U}]_{\overline{\gamma}\leftarrow\overline{\gamma}}$, hence is almost upper triangular.\\
\hlt{Structural Commentary}: Real triangularisation starts with an invariant line or plane.\\
\hlt{Key Insight}: Complex eigenvectors split into real and imaginary parts; those parts form the first real block.
\end{proof}
\end{theorem}

\begin{corollary} \label{cor:real-schur-block-form-orthonormal} \hlt{Orthonormal Real Schur Block Form}\\
Let $\dim_{\mathbb{R}}\mathcal{V}<\infty$ and let $T\in\LL(\mathcal{V})$. Then $\exists$ an orthonormal basis $\gamma$ of $\mathcal{V}$ s.t. $[T]_{\gamma\leftarrow\gamma}$ is in real Schur block form.
\begin{proof}
By Theorem \ref{thm:real-almost-upper-triangularization}, choose $\beta=(\mathbf{v}_1,\ldots,\mathbf{v}_n)$ with block sizes $d_j\in\{1,2\}$.\\
Put $s_k=d_1+\cdots+d_k$ and $\mathcal{U}_k=\operatorname{span}(\mathbf{v}_1,\ldots,\mathbf{v}_{s_k})$. Then $T(\mathcal{U}_k)\subseteq\mathcal{U}_k$ for all $k$.\\
Apply Gram--Schmidt to $\beta$ and obtain an orthonormal basis $\gamma=(\mathbf{e}_1,\ldots,\mathbf{e}_n)$ with $\operatorname{span}(\mathbf{e}_1,\ldots,\mathbf{e}_{s_k})=\mathcal{U}_k$, $1\leq k\leq m$. Thus $[T]_{\gamma\leftarrow\gamma}$ is block upper triangular for the same $d_j\in\{1,2\}$, hence is in real Schur block form.\\
\hlt{Structural Commentary}: Gram--Schmidt preserves the invariant block flag.\\
\hlt{Key Insight}: Orthonormalization changes blocks, not the flag.
\end{proof}
\end{corollary}

\begin{proposition} \label{prop:upper-triangular-invertibility-criterion} \hlt{Invertibility from an Upper-Triangular Matrix}\\
Suppose $T\in\LL(\mathcal{V})$ has an upper-triangular matrix with diagonal entries $\lambda_1,\ldots,\lambda_n$ with respect to some basis of $\mathcal{V}$. Then $T$ is invertible if and only if $\lambda_j\neq0$ for every $j$.
\begin{proof}
Let $\beta=(\mathbf{v}_1,\ldots,\mathbf{v}_n)$ be such a basis.
\begin{enumerate}[label=\roman*.]
\setlength{\itemsep}{0pt}
\item[$(\Leftarrow)$]
If every $\lambda_j\neq0$, then upper-triangular matrix $[T]_{\beta\leftarrow\beta}$ is invertible by Theorem \ref{thm:inverse-upper-triangular}, so $T$ is invertible.
\item[$(\Rightarrow)$]
Suppose $T$ is invertible. If $\lambda_j=0$ for some $j$, set $\mathcal{U}_j=\operatorname{span}(\mathbf{v}_1,\ldots,\mathbf{v}_j)$ and $\mathcal{U}_0=\{\mathbf{0}\}$. Upper triangularity gives $T(\mathcal{U}_j)\subseteq\mathcal{U}_{j-1}$ because the coefficient of $\mathbf{v}_j$ in $T\mathbf{v}_j$ is $\lambda_j=0$ and $T\mathbf{v}_i\in\mathcal{U}_{j-1}$ for $i<j$.\\
Since $\dim(\mathcal{U}_j)>\dim(\mathcal{U}_{j-1})$, the map $\mathcal{U}_j\to\mathcal{U}_{j-1}$ induced by $T$ is not injective. Thus there is a nonzero $\mathbf{u}\in\mathcal{U}_j$ with $T\mathbf{u}=\mathbf{0}$, contradicting invertibility. Hence every $\lambda_j\neq0$.
\end{enumerate}
\hlt{Structural Commentary}: Diagonal entries control each flag layer.\\
\hlt{Key Insight}: A zero diagonal entry forces a kernel vector.
\end{proof}
\end{proposition}

\begin{proposition} \label{prop:eigenvalues-of-upper-triangular-operator} \hlt{Eigenvalues from an Upper-Triangular Matrix}\\
Suppose $T\in\LL(\mathcal{V})$ has an upper-triangular matrix with diagonal entries $\lambda_1,\ldots,\lambda_n$ with respect to some basis of $\mathcal{V}$. Then the eigenvalues of $T$ are precisely $\lambda_1,\ldots,\lambda_n$.
\begin{proof}
For $\lambda\in\mathbb{F}$, the matrix of $T-\lambda I_{\mathcal{V}}$ in the same basis is upper triangular with diagonal entries $\lambda_1-\lambda,\ldots,\lambda_n-\lambda$. By Proposition \ref{prop:upper-triangular-invertibility-criterion}, $T-\lambda I_{\mathcal{V}}$ is not invertible if and only if $\lambda=\lambda_j$ for some $j$. By Theorem \ref{thm:equivalent-eigenvalue-conditions}, this is equivalent to $\lambda$ being an eigenvalue of $T$.\\
\hlt{Structural Commentary}: Scalar shifts preserve triangular form.\\
\hlt{Key Insight}: Eigenvalues occur exactly when one shifted diagonal entry becomes zero.
\end{proof}
\end{proposition}

\begin{perspective} \label{per:upper-triangular-viewpoint} \hlt{Why Triangular Form Matters}\\
Upper-triangular form is the first useful simplification beyond finding one eigenvector. It preserves enough of the operator to read eigenvalues from the diagonal while still allowing nontrivial coupling above the diagonal.
\end{perspective}
